\renewcommand\LTcaptype{quadro}
\begin{scriptsize}
\begin{longtable}[c]{|p{5cm}|l|l|l|l|l|l|l|}

\caption{Plano de tarefas do TCC 1}
\addcontentsline{loq}{quadro}{\protect\numberline{\thequadro}Plano de tarefas do TCC 1}
\label{tab:TCC1}\\
\hline
\multicolumn{1}{|c|}{Tarefas} &
  \multicolumn{1}{c|}{Mar/2025} &
  \multicolumn{1}{c|}{Abr/2025} &
  \multicolumn{1}{c|}{Mai/2025} &
  \multicolumn{1}{c|}{Jun/2025} &
  \multicolumn{1}{c|}{Jul/2025} &
  \multicolumn{1}{c|}{Ago/2025} &
  \multicolumn{1}{c|}{Set/2025} \\ \hline
\endhead
%
Realizar o levantamento bibliográfico acerca do tema "gasto energético em deficientes físicos que fazem uso de cadeira de rodas" &
  X &  &  &  &  &  &  \\ \hline
Identificar os principais fatores e parâmetros utilizados no cálculo de caloria &  & X &  &  &  &  &  \\ \hline
Mapear ferramentas e ou dispositivos voltados a identificação de métricas esportivas para deficientes &  & X & X &  &  &  &  \\ \hline
Selecionar métricas e/ou dispositivos que possam auxiliar na coleta de dados corporais de cadeirantes voltados ao gasto energético &  &  & X & X &  &  &  \\ \hline
Propor uma solução de monitoramento para pessoas cadeirantes &  &  &  & X & X &  &  \\ \hline
Elaboração e apresentação dos slides para o(a) orientador(a) &  &  &  &  &  & X &  \\ \hline
Apresentação do trabalho para a banca examinadora &  &  &  &  &  & X & X \\ \hline
Limite para a entrega dos ajustes após avaliação da banca examinadora &  &  &  &  &  &  & X \\ \hline
\caption*{Fonte: Elaborado pelo autor.}
\end{longtable}
\end{scriptsize}
\renewcommand\LTcaptype{quadro}
